\chapter{Velocity Composition, Light Cones, and Four-Vectors}

\section{Two Frames on One Railroad Track}

The Lorentz transformation derived previously concerned a deliberately restricted geometry: every relative motion lay along one axis. Picture a long railroad track. One set of observers remains at the station; another set rides in a train moving along the track. The transverse directions are present, but they play no role until they are restored later.

Each frame is an operational coordinate system. Meter sticks at rest in the frame define spatial coordinates, and synchronized clocks at rest in the frame define time. The postulate that fixes their relationship is that every inertial observer measures the same speed of light.

Most calculations will use \(c=1\). Time may be measured in seconds and distance in light-seconds, or in years and light-years. Velocities then become dimensionless fractions of light speed. When ordinary units are needed, factors of \(c\) are restored uniquely by dimensional consistency; for example, \(v^2\) in a dimensionless Lorentz factor becomes \(v^2/c^2\).

\subsection{Direct and inverse transformations}

Let a flash occur somewhere in or near the train. It is one event, independent of who describes it. The station observer assigns \((x,t)\), while the train observer assigns \((x',t')\). In units \(c=1\),
\begin{align}
x'&=\gamma_v(x-vt),&
t'&=\gamma_v(t-vx),
\label{eq:sr2-direct}\\
\gamma_v&=\frac{1}{\sqrt{1-v^2}}.
\end{align}
Newton would recognize the displacement \(x-vt\); the factor \(\gamma_v\) and the mixing of \(x\) into \(t'\) are the relativistic changes.

\begin{figure}[htbp]
\centering
\includegraphics[width=0.90\textwidth]{lecture_02_figure_01.png}
\caption{The direct Lorentz transformation relating the station coordinates \((x,t)\) to the train coordinates \((x',t')\).}
\lectureframe{00:07:52}
\end{figure}

Reciprocity gives the inverse by reversing the sign of the relative velocity:
\begin{align}
x&=\gamma_v(x'+vt'),&
t&=\gamma_v(t'+vx').
\label{eq:sr2-inverse}
\end{align}
The station is not absolutely at rest. It is merely the frame called unprimed, and the train sees it moving at \(-v\).

\begin{figure}[htbp]
\centering
\includegraphics[width=0.90\textwidth]{lecture_02_figure_02.png}
\caption{The inverse map has the same Lorentz factor and the opposite sign of relative velocity.}
\lectureframe{00:09:15}
\end{figure}

For a boost along \(x\), the transverse coordinates do not change:
\begin{equation}
y'=y,
\qquad
z'=z.
\end{equation}

The relative speed can be read from the space equation without studying its denominator. The train origin is defined by \(x'=0\), and therefore
\begin{equation}
x'=0
\quad\Longleftrightarrow\quad
x-vt=0
\quad\Longleftrightarrow\quad
x=vt.
\label{eq:sr2-origin-worldline}
\end{equation}
The coefficient of \(t\) in this worldline is the train's speed in station coordinates.

\section{Three Frames and the Composition of Boosts}

Place a small kiddie car inside the train. The train moves at speed \(v\) relative to the station, while the car moves down the aisle at speed \(u\) relative to the train passenger. Call the speed of the car relative to the station \(w\). Newtonian intuition suggests \(w=u+v\), but the Lorentz transformations must decide.

There are now three coordinate systems:
\begin{equation}
(x,t),\qquad (x',t'),\qquad (x'',t'').
\end{equation}
The first belongs to the station, the second to the train, and the third to the car. Between the train and the car,
\begin{align}
x''&=\gamma_u(x'-ut'),&
t''&=\gamma_u(t'-ux'),
\label{eq:sr2-second-boost}\\
\gamma_u&=\frac{1}{\sqrt{1-u^2}}.
\end{align}

\begin{figure}[htbp]
\centering
\includegraphics[width=0.90\textwidth]{lecture_02_figure_03.png}
\caption{The direct station-to-train formulas are followed by a second Lorentz transformation from primed to double-primed coordinates.}
\lectureframe{00:14:55}
\end{figure}

\subsection{Substitute one boost into the other}

The spatial equation is enough to find \(w\). Substitute \eqref{eq:sr2-direct} into the first equation of \eqref{eq:sr2-second-boost}:
\begin{align}
x''
&=\gamma_u(x'-ut')\\
&=\gamma_u\gamma_v\left[(x-vt)-u(t-vx)\right]\\
&=\gamma_u\gamma_v\left[(1+uv)x-(u+v)t\right].
\label{eq:sr2-composed-space}
\end{align}
The two minus signs in the \(uvx\) term make it positive; the two time terms combine as \(-(u+v)t\).

\begin{figure}[htbp]
\centering
\includegraphics[width=0.91\textwidth]{lecture_02_figure_04.png}
\caption{The nested substitution before collecting the \(x\)- and \(t\)-terms in the composed boost.}
\lectureframe{00:16:30}
\end{figure}

As before, the denominator is irrelevant when locating the moving origin. The kiddie car is at rest at \(x''=0\), so the bracket in \eqref{eq:sr2-composed-space} must vanish:
\begin{equation}
(1+uv)x-(u+v)t=0.
\end{equation}
Its station-frame worldline is therefore
\begin{equation}
x=\frac{u+v}{1+uv}\,t,
\end{equation}
and its velocity is
\begin{equation}
\boxed{w=\frac{u+v}{1+uv}}.
\label{eq:sr2-velocity-add-c1}
\end{equation}

\begin{figure}[htbp]
\centering
\includegraphics[width=0.91\textwidth]{lecture_02_figure_05.png}
\caption{The collected numerator and the worldline \(x=[(u+v)/(1+uv)]t\) reveal the composed velocity.}
\lectureframe{00:19:30}
\end{figure}

The time equation composes consistently as well. The complete station-to-car transformation has the Lorentz form
\begin{align}
x''&=\gamma_w(x-wt),&
t''&=\gamma_w(t-wx),
\end{align}
where
\begin{equation}
\gamma_w=\frac{1}{\sqrt{1-w^2}}.
\end{equation}
Indeed,
\begin{equation}
\gamma_w=\gamma_u\gamma_v(1+uv),
\end{equation}
which follows by inserting \eqref{eq:sr2-velocity-add-c1} and simplifying. Two collinear boosts therefore compose to another collinear boost.

\subsection{Restoring dimensions}

\begin{classroomqa}
\audiencequestion{Why is the denominator \(1+uv\) dimensionally incorrect in ordinary units?}
\lecturerresponse{It is correct only because \(c=1\) makes every velocity a dimensionless fraction of light speed. In conventional units the product must appear as \(uv/c^2\), which is dimensionless. \lecturetimestamp{00:21:11}}
\end{classroomqa}

The dimensional velocity-addition law is
\begin{equation}
\boxed{
w=\frac{u+v}{1+\dfrac{uv}{c^2}}
}.
\label{eq:sr2-velocity-add}
\end{equation}

\begin{figure}[htbp]
\centering
\includegraphics[width=0.90\textwidth]{lecture_02_figure_06.png}
\caption{Dimensional consistency restores \(c^2\) beneath the product \(uv\).}
\lectureframe{00:22:06}
\end{figure}

For \(\lvert u\rvert,\lvert v\rvert\ll c\),
\begin{equation}
\frac{uv}{c^2}\ll 1,
\end{equation}
so
\begin{equation}
w\approx u+v.
\end{equation}
Newtonian addition is the low-speed approximation rather than a separate rule.

\subsection{Two numerical regimes}

First take
\begin{equation}
u=v=0.01c.
\end{equation}
In units \(c=1\),
\begin{equation}
w=\frac{0.02}{1.0001}\approx 0.019998.
\end{equation}
One percent of light speed is already about \(3\times10^6\ {\rm m/s}\), yet the relativistic correction is only one part in ten thousand. The result is slightly smaller than the Newtonian \(0.02c\).

\begin{figure}[htbp]
\centering
\includegraphics[width=0.90\textwidth]{lecture_02_figure_07.png}
\caption{At \(u=v=0.01c\), the denominator differs from one by only \(10^{-4}\).}
\lectureframe{00:24:23}
\end{figure}

Now take
\begin{equation}
u=v=0.9c.
\end{equation}
Newton would predict \(1.8c\). Relativity gives
\begin{equation}
\frac{w}{c}
=\frac{0.9+0.9}{1+0.9^2}
=\frac{1.8}{1.81}
\approx 0.9945.
\label{eq:sr2-fast-example}
\end{equation}

\begin{figure}[htbp]
\centering
\includegraphics[width=0.90\textwidth]{lecture_02_figure_08.png}
\caption{Two speeds of \(0.9c\) compose to \(1.8/1.81\), still below light speed.}
\lectureframe{00:25:48}
\end{figure}

Successive subluminal boosts cannot cross \(c\). If \(u<c\) and \(v<c\), then
\begin{equation}
1-\frac{w}{c}
=
\frac{\left(1-\frac{u}{c}\right)
      \left(1-\frac{v}{c}\right)}
     {1+\frac{uv}{c^2}}
>0.
\end{equation}
Each additional boost can bring the result closer to \(c\), but never beyond it. At the limiting value \(u=c\),
\begin{equation}
w=\frac{c+v}{1+v/c}=c.
\end{equation}
The speed of light is a fixed point of the composition law.

\section{Coordinate Measurement Is Not Visual Appearance}

\begin{classroomqa}
\audiencequestion{If the train has glass walls, could the station observer simply look at the passenger and kiddie car together and read their motions directly?}
\lecturerresponse{The Lorentz transformation concerns coordinate measurements made with distributed rods and synchronized clocks. A visual image also contains the unequal travel times of light from different events to the observer. Glass changes the view, not the coordinate relation, and treating one image as simultaneous data would reintroduce the very simultaneity assumption that relativity corrected. \lecturetimestamp{00:27:53}}
\end{classroomqa}

An event has a location and a time in each frame whether or not anyone happens to see it at that moment. Optical appearance is a legitimate subject, but it requires a second calculation tracing light from each event to the eye. The transformation laws themselves are universal and independent of the train's walls.

\section{The Proper Interval in Three Spatial Dimensions}

Return to two events, initially one at the origin and the other at \((x,t)\). In one spatial dimension,
\begin{equation}
\tau^2=t^2-x^2
\label{eq:sr2-interval-1d}
\end{equation}
in units \(c=1\). Every inertial frame agrees:
\begin{equation}
\tau^2
=t^2-x^2
=t'^2-x'^2
=t''^2-x''^2.
\end{equation}
The coordinates change; the proper interval does not.

For timelike separation, a clock can move between the events. If it reads noon at the first event, its elapsed reading at the second is the proper time \(\tau\). That statement is independent of which coordinate frame is used to calculate the interval.

\begin{figure}[htbp]
\centering
\includegraphics[width=0.90\textwidth]{lecture_02_figure_09.png}
\caption{The invariant proper-time relation \(\tau^2=t^2-x^2\) is recalled before restoring the transverse directions.}
\lectureframe{00:30:18}
\end{figure}

\subsection{Spatial rotations restore the transverse directions}

Consider an ordinary rotation of spatial axes before returning to boosts. A point not lying on the \(x\)-axis has squared spatial distance
\begin{equation}
r^2=x^2+y^2+z^2,
\end{equation}
which is invariant under rotations by the ordinary Pythagorean theorem. The \(3+1\)-dimensional spacetime interval is consequently
\begin{equation}
\boxed{
\tau^2=t^2-x^2-y^2-z^2
}.
\label{eq:sr2-interval-3d}
\end{equation}
With \(c\) restored,
\begin{equation}
c^2\tau^2
=c^2t^2-x^2-y^2-z^2.
\end{equation}

\begin{figure}[htbp]
\centering
\includegraphics[width=0.90\textwidth]{lecture_02_figure_10.png}
\caption{Pythagorean spatial distance extends the invariant from \(t^2-x^2\) to \(t^2-x^2-y^2-z^2\).}
\lectureframe{00:34:23}
\end{figure}

Any combination of spatial rotations and Lorentz boosts preserves \eqref{eq:sr2-interval-3d}. Ordinary coordinate axes disagree about components but agree about Euclidean distance; inertial spacetime frames disagree about space and time components but agree about this difference of squares. That invariant is the central geometric fact of special relativity.

\begin{classroomqa}
\audiencequestion{If the passenger replaces the kiddie car by a forward flashlight beam, do the train and station observers obtain the same light speed?}
\lecturerresponse{Set \(u=c\) in \eqref{eq:sr2-velocity-add}. Then \(w=(c+v)/(1+v/c)=c\). In \(c=1\) units the ratio is simply \((1+v)/(1+v)=1\). Both frames measure the same speed, exactly as required. \lecturetimestamp{00:35:53}}
\end{classroomqa}

The earlier glass-wall concern is therefore not a failure of invariant light speed. It is a confusion between a coordinate measurement and a delayed visual image, ultimately a confusion about simultaneity.

\section{Null Rays and the Light Cone}

In one space dimension, right- and left-moving light rays satisfy
\begin{equation}
x=\pm t.
\end{equation}
Their interval vanishes:
\begin{equation}
t^2-x^2=0.
\end{equation}
Unlike Euclidean distance, a zero spacetime interval does not require coincident events. It means that a light signal can connect them.

With all spatial directions restored, a light ray obeys
\begin{equation}
t^2-x^2-y^2-z^2=0,
\label{eq:sr2-null}
\end{equation}
or
\begin{equation}
r=t
\end{equation}
for future-directed light in \(c=1\) units.

\begin{figure}[htbp]
\centering
\includegraphics[width=0.90\textwidth]{lecture_02_figure_11.png}
\caption{Light in any spatial direction follows a null trajectory satisfying \(t^2-(x^2+y^2+z^2)=0\).}
\lectureframe{00:39:28}
\end{figure}

In \(1+1\) dimensions the null set is a pair of \(45^\circ\) lines. With two spatial dimensions plus time it becomes a cone; three spatial dimensions cannot be displayed fully on a board, but the defining equation is unchanged.

The future light cone consists of events that light emitted at the origin can reach. The past light cone consists of events from which light can reach the origin. One is the other turned over in time.

\begin{figure}[htbp]
\centering
\includegraphics[width=0.90\textwidth]{lecture_02_figure_12.png}
\caption{The \(45^\circ\) null rays sweep out the future and past light cones in a higher-dimensional spacetime drawing.}
\lectureframe{00:40:47}
\end{figure}

\begin{classroomqa}
\audiencequestion{What is the corresponding invariant under Galilean transformations in Newtonian mechanics?}
\lecturerresponse{The time interval itself. Newtonian observers all agree on \(\Delta t\), even though they disagree on spatial displacement. Relativity relinquishes universal time and replaces it with the invariant spacetime interval. \lecturetimestamp{00:41:53}}
\end{classroomqa}

\section{From Spatial Vectors to Four-Vectors}

The vectors here are geometric vectors, not quantum state vectors. A displacement between two points in ordinary space has magnitude and direction. Its tail can be translated to the origin without changing the vector, and its three components may be denoted
\begin{equation}
x^i=(x^1,x^2,x^3)=(x,y,z),
\qquad i=1,2,3.
\end{equation}
The Latin index \(i\) runs only over spatial directions.

An event needs one additional component, its time. A spacetime position is therefore written
\begin{equation}
x^\mu=(x^0,x^1,x^2,x^3)=(t,x,y,z),
\qquad \mu=0,1,2,3.
\label{eq:sr2-four-position}
\end{equation}
By historical convention time is the zeroth component rather than the fourth:
\begin{equation}
x^0=t,\qquad x^1=x,\qquad x^2=y,\qquad x^3=z.
\end{equation}
Greek indices such as \(\mu\) run over spacetime; Latin indices such as \(i\) run over space.

\begin{figure}[htbp]
\centering
\includegraphics[width=0.90\textwidth]{lecture_02_figure_13.png}
\caption{Spatial components \(x^i\) are enlarged to the four-component spacetime position \(x^\mu=(x^0,x^1,x^2,x^3)\).}
\lectureframe{00:46:05}
\end{figure}

The invariant becomes
\begin{equation}
\tau^2
=(x^0)^2-(x^1)^2-(x^2)^2-(x^3)^2.
\end{equation}
The notation alone adds no new physics. Its value is that a Lorentz transformation acts linearly on all four components, just as a rotation acts linearly on a spatial vector.

\begin{classroomqa}
\audiencequestion{Since the transformation is linear, can it be represented by a matrix acting on a column vector?}
\lecturerresponse{Yes. For a boost along \(x\), the nontrivial block is \(\gamma_v\begin{pmatrix}1&-v\\-v&1\end{pmatrix}\) acting on \((t,x)^{\mathsf T}\), while \(y\) and \(z\) are unchanged. Matrix notation is the natural way to compose and organize Lorentz transformations. \lecturetimestamp{00:48:14}}
\end{classroomqa}

Explicitly,
\begin{equation}
\begin{pmatrix}
t'\\x'\\y'\\z'
\end{pmatrix}
=
\begin{pmatrix}
\gamma_v&-\gamma_v v&0&0\\
-\gamma_v v&\gamma_v&0&0\\
0&0&1&0\\
0&0&0&1
\end{pmatrix}
\begin{pmatrix}
t\\x\\y\\z
\end{pmatrix}.
\label{eq:sr2-boost-matrix}
\end{equation}

\section{Displacements and Four-Velocity}

Move from the position of one event to a short segment of a particle's worldline. The four-displacement is
\begin{equation}
\Delta x^\mu
=(\Delta t,\Delta x,\Delta y,\Delta z).
\end{equation}
When a displacement is called small, calculus is approaching: \(\Delta x^\mu\) will become \(dx^\mu\). For the moment, the finite notation makes the geometry transparent.

The proper-time increment along a timelike worldline satisfies
\begin{equation}
(\Delta\tau)^2
=(\Delta t)^2-(\Delta x)^2-(\Delta y)^2-(\Delta z)^2.
\label{eq:sr2-proper-increment}
\end{equation}
Ordinary velocity divides spatial displacement by coordinate time and has three components. Four-velocity instead divides the four-displacement by the invariant proper time:
\begin{equation}
u^\mu
=\frac{\Delta x^\mu}{\Delta\tau}
\quad\longrightarrow\quad
\boxed{
u^\mu=\frac{dx^\mu}{d\tau}
}.
\label{eq:sr2-four-velocity}
\end{equation}

\begin{figure}[htbp]
\centering
\includegraphics[width=0.90\textwidth]{lecture_02_figure_14.png}
\caption{A worldline segment defines \(\Delta x^\mu\), its invariant \(\Delta\tau\), and the four-velocity \(u^\mu=\Delta x^\mu/\Delta\tau\).}
\lectureframe{00:52:42}
\end{figure}

This object has four components and transforms as a four-vector. Its relation to ordinary three-velocity is the next question, deliberately left open here. Resolving it leads toward relativistic momentum, energy, acceleration, and ultimately the relativistic generalization of \(F=ma\).
